\documentclass[a4paper, 11pt]{article}

\usepackage[T1,T2A]{fontenc}
\usepackage[utf8]{inputenc}
\usepackage[english, russian]{babel}

\usepackage{comment}
\usepackage{fullpage}
\usepackage{amssymb}
\usepackage{amsmath}
\usepackage{amsthm}
\usepackage{enumitem}

% %\DeclareMathOperator{\ord}{ord}
\DeclareMathOperator{\chr}{char}
\DeclareMathOperator{\Mod}{mod}
\DeclareMathOperator{\Gal}{Gal}
\DeclareMathOperator{\N}{N}
\DeclareMathOperator{\Tr}{Tr}
\DeclareMathOperator{\SL}{SL}

\begin{document}

\begin{center} \Large \bf Листок 2-3\\ Тема 2.3. Корни из единицы. Круговой многочлен \end{center}

\noindent\textbf{Упражнения и задачи}
\begin{enumerate}[topsep=0pt]
    \item Пусть $k$ --- поле, $f=a_nx^n+\dots+a_1x+a_0 \in k[x]$, $f'=n a_n x^{n-1}+\dots+a_1 \in k[x]$ --- формальная производная $f$. Докажите, что для любых $f, g \in k[x]$ верны следующие свойства:
    \begin{itemize}[topsep=0pt]
        \item $(f+g)'=f'+g'$;
        \item $(fg)'=f'g+fg'$.
    \end{itemize}

    \item Пусть $k$ --- поле, $f \in k[x]$. Докажите, что $\alpha \in k$ --- кратный корень $f$ $\Leftrightarrow$ $f(\alpha) = f'(\alpha) = 0$.

    \item Пусть $f \in \mathbb{F}_{p^m}[x]$. Докажите, что $f\in \mathbb{F}_{p}[x] \Leftrightarrow (f(x))^p=f(x^p)$.

    \item Докажите, что если $\mathrm{char}\, k = p \nmid n$, то для кругового многочлена верно\\ $\Phi_n(x) = \prod\limits_{d|n}(x^d-1)^{\mu(n/d)}$.
    
    \item Докажите, что $\mathbb{F}_q$ есть $(q-1)$-е круговое поле над любым своим подполем. %[LN, p63 Lemma 2.49]

    \item Пусть $\alpha \in \mathbb{F}_q$, $n \in \mathbb{Z}_{\geqslant 0}$. Докажите, что $x^q-x+\alpha\ |\ x^{q^n}-x+n\alpha$. %[LN p71 ex 2.19]

    \item Пусть $f \in \mathbb{F}_q[x]$, $q=p^n$. Докажите, что $f'(x)=0$ $\Leftrightarrow$ $f=g^p$ для некоторого $g \in \mathbb{F}_q[x]$. %[LN p71 ex 2.23]

    \item Пусть $\mathbb{F}_{q^m}/\mathbb{F}_q$ --- расширение конечных полей, $\alpha \in \mathbb{F}_{q^m}^*$. Докажите, что сопряжённые элементы $\alpha, \alpha^q, \alpha^{q^2}, \dots, \alpha^{q^{m-1}}$ имеют один и тот же порядок в мультипликативной группе $\mathbb{F}_{q^m}^*$.   %[LN, p50 th2.18]

    \item Пусть $f\in\mathbb{F}_q[x]$ --- неприводимый многочлен степени $d$, $m\geqslant 1$. Докажите, что $f$ неприводим над $\mathbb{F}_{q^m}$ $\Leftrightarrow$ $(d,m)=1$. Более того, если $(d,m)=l$, то над $\mathbb{F}_{q^m}$ многочлен $f$ раскладывается в произведение $l$ неприводимых многочленов степени $d/l$. %[LN Th 3.46]

    \item Пусть $f\in \mathbb{F}_q[x]$, $q=p^n$, $\deg f = m \geqslant 1$, $f(0) \neq 0$. Докажите, что $\exists e \in \mathbb{Z}_+$, $e\leqslant q^m-1$, такое, что $f(x)\ |\ x^e-1$. Наименьшее такое $e$ называется порядком многочлена $f(x)$ в $\mathbb{F}_q[x]$. Докажите также следующие свойства:
    \begin{itemize}[noitemsep,topsep=0pt]
        \item Пусть $f$ --- неприводимый, тогда порядок $f$ равен порядку $\alpha \in \mathbb{F}_{q^m}^*$, где $\alpha$ --- корень $f$;
        \item Пусть $f$ --- неприводимый, тогда порядок $f$ делит $q^m-1$;
        \item Пусть $c\in \mathbb{Z}_+$, $e$ --- порядок $f$, тогда $f(x)\ |\ x^c-1$ $\Leftrightarrow$ $e|c$;
        \item Пусть $e_1, e_2\in \mathbb{Z}_+$, тогда наибольший общий делитель многочленов $x^{e_1}-1$, $x^{e_2}-1$ в $\mathbb{F}_q[x]$ равен $x^d-1$, где $d=(e_1,e_2)$.
    \end{itemize} %[LN, p77-78 3.1-3.7]

\end{enumerate}

% \noindent\textbf{SageMath}
% \begin{itemize}[topsep=0pt]
%     \item Исследуйте основные функции SageMath связанные с работой в кольцах многочленов над конечными полями:
%     \begin{itemize}[noitemsep,topsep=0pt]
%         \item Кольцо многочленов: \texttt{PolynomialRing()};
%         \item Неприводимость многочлена: \texttt{is\_irreducible()};
%         \item Разложение многочлена на множители: \texttt{factor()};
%         \item Корни многочлена: \texttt{roots()};
%         \item Круговой многочлен: \texttt{cyclotomic\_polynomial()}.
%      \end{itemize}
        
% \end{itemize}

\noindent\textbf{Темы для самостоятельного изучения}
\begin{itemize}[topsep=0pt]
    \item Изучите «элементарные» доказательства неприводимости кругового многочлена $\Phi_n(x)$ в $\mathbb{Z}[x]$.
    (см. например\\ \texttt{https://www.lehigh.edu/~shw2/c-poly/several\_proofs.pdf})
    \item Теорема Веддербёрна: Всякое конечное тело (кольцо с единицей, в котором каждый ненулевой элемент обратим) является полем. ([LN], [The Book]).
\end{itemize}

\end{document}
